\documentclass[letterpaper]{article}
\usepackage[margin=1in]{geometry}
\usepackage{amsmath,amssymb}
\usepackage{graphicx}
\usepackage{booktabs}
\usepackage{hyperref}
\usepackage{xcolor}
\usepackage{caption}
\usepackage{subcaption}
\usepackage{algorithm}
\usepackage{algorithmic}
\usepackage{siunitx}
\usepackage{float}       % [H] placement specifier
\usepackage{placeins}    % \FloatBarrier

% Citation style — update to venue format (e.g., IEEEtran for RA-L, acmart for CoRL)
\bibliographystyle{plain}

% ── Colour macros ─────────────────────────────────────────────────────────────
\definecolor{ProvenGreen}{HTML}{22c55e}
\definecolor{DangerRed}{HTML}{ef4444}
\definecolor{CbfAmber}{HTML}{f59e0b}
\newcommand{\proven}[1]{\textcolor{ProvenGreen}{\textbf{#1}}}
\newcommand{\danger}[1]{\textcolor{DangerRed}{\textbf{#1}}}
\newcommand{\cbf}[1]{\textcolor{CbfAmber}{\textbf{#1}}}

% ── Maths shorthands ─────────────────────────────────────────────────────────
\newcommand{\hcbf}{h_{\text{CBF}}}
\newcommand{\dsafe}{d_{\text{safe}}}
\newcommand{\dsurf}{d_{\text{surf}}}
\newcommand{\vdes}{\mathbf{v}_{\text{des}}}
\newcommand{\vcmd}{\mathbf{v}_{\text{cmd}}}

% ─────────────────────────────────────────────────────────────────────────────
\title{\textbf{FleetSafe-VisualNav: Paradigm-Selective Command-Layer Safety
       for Visual Navigation via CBF Intervention}\\[4pt]
       \large Architecture-Agnostic Command-Layer Safety Filtering Without Policy Modification}
\author{F. Van Laarhoven \and [Co-authors TBD]\\
  \small Newcastle University, UK\\
  \small \texttt{F.Van-Laarhoven2@newcastle.ac.uk} \quad
  ORCID: 0009-0006-8931-0364}
\date{May 2026 — DRAFT}
% ─────────────────────────────────────────────────────────────────────────────

\begin{document}
\maketitle
\begin{abstract}
Visual navigation policies for mobile robots achieve impressive task performance
but present distinct failure structures depending on their underlying navigation
paradigm. We present \textbf{FleetSafe}, a command-layer Control Barrier Function
(CBF) filter that intervenes on velocity commands issued by any black-box
navigation policy without modifying the underlying model weights.
We evaluate three policies spanning the major paradigms of learned visual
navigation---GNM (goal-conditioned generalist), ViNT (Transformer-based
foundation model), and NoMaD (diffusion-based exploration)---in two complementary
simulation backends: MuJoCo (visible procedural obstacles, 50 seeds) and
Isaac Sim (invisible map-registered hazard cylinders, 50 seeds).

Our experiments reveal that invisible hazard vulnerability is
\emph{navigation-paradigm-dependent}: goal-directed policies (GNM, ViNT)
commit to obstacle-crossing trajectories toward their targets and collide at
100\% RAW collision rate, while the diffusion-based NoMaD avoids the same hazard
naturally (0\% RAW, min clearance\,=\,1.5\,m). FleetSafe correctly intervenes
for GNM/ViNT (CBF IR\,=\,49.7\%/37.9\%, reducing collision to 0\%) and
remains idle for NoMaD (IR\,=\,0\%, do-no-harm confirmed under tested
conditions), demonstrating \emph{selective} intervention rather than
over-conservative suppression of all commands.

In MuJoCo, unsafe behaviour is model-specific (ViNT's aggressive forward prior,
100\%\,$\to$\,0\%); GNM and NoMaD exhibit conservative behaviour throughout.
The same CBF-QP filter, with identical hyperparameters, handles all cases
without access to policy internals. The filter operates in $O(n_{\text{obs}})$
time with a closed-form single-constraint solution, achieving
\SI{20}{\milli\second} mean inference latency (well below \SI{100}{\milli\second}
real-time threshold). All results carry $n\,{=}\,50$ seeds and PROVEN-gate
certification; evidence artefacts are publicly released.
\end{abstract}

% ─────────────────────────────────────────────────────────────────────────────
\section{Introduction}
\label{sec:intro}

Foundation model visual navigation policies have demonstrated remarkable
generalisation across environments \cite{shah2023gnm,sridhar2024nomad,
majumdar2023vint}. However, their black-box nature and training-distribution
sensitivity make it difficult to provide formal safety guarantees in
safety-critical deployments such as hospital robotics.

We identify two distinct failure mechanisms that a command-layer safety
filter must handle:
\begin{enumerate}
  \item \textbf{Model-specific aggressive priors}: Certain VLA backbones
    (notably ViNT) exhibit an aggressive forward bias that causes
    \danger{100\%} collision rates in cluttered hospital corridors even
    when obstacles are visually prominent. This is a model-\emph{dependent}
    failure---different backbones behave differently in identical scenes.
  \item \textbf{Map-registered invisible hazards}: In environments where
    hazards are geometrically real but camera-invisible (e.g., glass
    partitions, transparent doors, registered-but-not-visible obstacles),
    \emph{goal-directed} navigation backbones (GNM, ViNT) fail at
    \danger{100\%} collision because they navigate toward their goal
    through an obstacle the camera cannot see.
    Notably, the diffusion-based NoMaD avoids naturally (0\% RAW collision,
    min\_dist\,=\,1.5\,m) because its exploration prior does not commit to
    a fixed collision path. This reveals a \emph{navigation-paradigm
    dependency}: goal-directed VLAs are uniquely vulnerable to invisible
    map-registered hazards.
\end{enumerate}

FleetSafe addresses both failure types with the same command-layer CBF-QP filter:
the filter operates on the robot's on-board map rather than camera input,
intervening whenever the commanded velocity would violate a safety constraint.
The \emph{same} filter, with identical hyperparameters, produces:
\begin{itemize}
  \item \textbf{MuJoCo}: CBF idle for GNM/NoMaD (conservative) but
    \cbf{actively} corrects ViNT (IR\,=\,69.7\%), reducing collision from
    \danger{100\%} to \proven{0\%}.
  \item \textbf{Isaac Sim}: CBF intervenes for GNM (IR\,=\,49.7\%) and ViNT
    (IR\,=\,37.9\%); idle for NoMaD which naturally avoids
    (do-no-harm confirmed, IR\,=\,0\%). All collision rates: \proven{0\%}.
\end{itemize}
This demonstrates \emph{architecture-agnostic command-layer safety filtering}:
the mechanism operates at the velocity-command interface and requires no knowledge
of, nor access to, the VLA policy internals.

\paragraph{Contributions.}
\begin{enumerate}
  \item \textbf{Command-layer CBF-QP wrapper} for visual navigation policies,
    operating at the velocity interface without modifying policy weights or
    internal representations.
  \item \textbf{Cross-paradigm validation} across three fundamentally different
    navigation approaches: goal-conditioned generalist (GNM), Transformer
    foundation model (ViNT), and diffusion-based exploration (NoMaD).
    This span makes the architecture-agnostic intervention claim substantially
    more meaningful than evaluating minor variants of a single model class.
  \item \textbf{Navigation-paradigm-dependent failure analysis}: evidence that
    invisible map-registered hazard vulnerability depends on the navigation
    paradigm (goal-directed commitment vs.\ diffusion exploration), not merely
    on model capacity.
  \item \textbf{Selective intervention evidence}: FleetSafe intervenes only
    where required and remains idle (IR\,=\,0\%) for naturally safe policies,
    directly rebutting the over-conservatism criticism of safety filters.
  \item \textbf{Reproducible dual-backend benchmark} with $n\,{=}\,50$ seeds per
    condition, automated PROVEN gates, and full evidence artefacts.
\end{enumerate}

% ─────────────────────────────────────────────────────────────────────────────
\section{Background and Related Work}
\label{sec:related}

\subsection{Visual Navigation Foundation Models}
GNM \cite{shah2023gnm}, ViNT \cite{majumdar2023vint}, and NoMaD
\cite{sridhar2024nomad} are transformer-based navigation policies
trained on large real-world topological datasets encompassing millions
of navigation frames across diverse indoor and outdoor environments.
GNM uses a metric-map-conditioned goal-following objective with a
conservative velocity prior. ViNT augments this with a diffusion-based
planner that produces aggressive exploration behaviour. NoMaD adopts a
goal-masked diffusion approach that decouples exploration from directed
navigation. All three operate from onboard monocular RGB images alone,
making them blind to camera-invisible hazards.

\subsection{Safety in Robot Learning}
Safe reinforcement learning methods modify the training objective to
incorporate safety constraints \cite{garcia2015survey}, but are
inapplicable to frozen pre-trained VLA policies. Shielding approaches
\cite{alshiekh2018safe} apply a model-free correctional policy, which
still requires retraining. Constrained policy search \cite{achiam2017cpo}
modifies the policy gradient, again requiring gradient access.
Runtime monitoring \cite{falcone2011monitoring} detects but does not
prevent unsafe actions.

\subsection{Control Barrier Functions}
CBFs provide formal safety guarantees via set invariance
\cite{ames2019cbf}. The barrier function $h : \mathcal{X} \to \mathbb{R}$
satisfies $\{x : h(x) \geq 0\} = \mathcal{C}$ (safe set). The CBF
condition requires $\dot{h}(x) \geq -\alpha(h(x))$ for all $x \in \mathcal{C}$.
Ames et al.\ \cite{ames2017cbf} showed this can be enforced via a QP
with closed-form solutions in the single-constraint case. Recent work has
extended CBFs to handle control-affine dynamics \cite{xiao2021high},
input constraints \cite{rauscher2016cbf}, and robot navigation
\cite{wang2017safe}.

\subsection{Command-Layer Safety Filtering}
Command-layer (velocity-level) filters intervene on policy outputs
without accessing policy internals, making them architecture-agnostic.
Robey et al.\ \cite{robey2021cbf} and Emam et al.\ \cite{emam2022cbf}
demonstrated this for robotic manipulation and multi-agent navigation
respectively. Our work extends this paradigm to hospital-domain VLA
navigation, contributing (1) evidence across three distinct VLA
architectures, (2) a dual-physics-engine evaluation that separates
model-specific failure from navigation-paradigm-specific failure,
and (3) the surface-distance CBF formulation correcting the center-to-center
approximation for finite-radius obstacles.

\subsection{Hospital Robot Safety}
Hospital navigation raises unique safety requirements beyond standard
indoor navigation: narrow corridors, human traffic, glass partitions,
and map-registered but camera-invisible hazards (e.g., sterile zone
markers, temporary blockades). Prior work \cite{ozkil2009service,
cardinaux2011video} focused on SLAM and human detection; we focus on
command-layer safety filtering applicable when perception fails.

% ─────────────────────────────────────────────────────────────────────────────
\section{Method}
\label{sec:method}

\subsection{Problem Formulation}
Consider a ground robot with state $(x, y, \theta)$ navigating a hospital
environment. A VLA policy $\pi$ outputs desired velocity commands
$\vdes = (\dot{x}_d, \dot{y}_d)$. The FleetSafe filter maps
$\vdes \mapsto \vcmd$ via a QP that minimises deviation from the desired
command while enforcing safety constraints.

\subsection{Surface-Distance CBF}
For a cylindrical obstacle centred at $\mathbf{o}_i \in \mathbb{R}^2$
with radius $r_i$, we define the \emph{surface distance}:
\begin{equation}
  \dsurf(\mathbf{p}, i) = \|\mathbf{p} - \mathbf{o}_i\|_2 - r_i
\end{equation}
and the barrier function:
\begin{equation}
  \hcbf(\mathbf{p}, i) = \dsurf(\mathbf{p}, i)^2 - \dsafe^2
  \label{eq:cbf}
\end{equation}
The safe set is $\mathcal{C} = \{\mathbf{p} : \hcbf(\mathbf{p}, i) \geq 0\
\forall i\}$, i.e., the robot surface is at least $\dsafe$ from every
obstacle surface. For $r_i = 0$ (point obstacles, MuJoCo baseline),
\eqref{eq:cbf} reduces to the standard center-to-center formula.

\subsection{CBF-QP}
At each control step, the filter solves:
\begin{align}
  \vcmd^* &= \arg\min_{\mathbf{v}} \|\mathbf{v} - \vdes\|^2 \label{eq:qp}\\
  \text{s.t.} \quad & \nabla_{\mathbf{p}} \hcbf^{(i)} \cdot \mathbf{v} \geq -\alpha \hcbf^{(i)}(\mathbf{p})
    \quad \forall i \nonumber
\end{align}
with $\alpha > 0$ the class-$\mathcal{K}$ gain.

\paragraph{Closed-form solution.}
When a single constraint $\mathbf{a}^\top \mathbf{v} \geq b$ is active, the
QP~\eqref{eq:qp} has a closed-form solution:
\begin{equation}
  \vcmd = \vdes + \max\!\left(0,\ \frac{b - \mathbf{a}^\top\vdes}{\|\mathbf{a}\|^2}\right)\mathbf{a}
\end{equation}
This requires $O(n_{\text{obs}})$ time (one dot product per obstacle), making
it suitable for real-time robot control at \SI{10}{\hertz}. When multiple
constraints are simultaneously active, we solve the full QP (convex, $2d$
variables, $n_{\text{obs}}$ inequality constraints); in practice, at most
2--3 obstacles are simultaneously relevant, so the solver completes in
$<\SI{1}{\milli\second}$.

Algorithm~\ref{alg:cbf} presents the per-step filter implementation.

\begin{algorithm}[htbp]
\caption{FleetSafe CBF-QP Filter (per control step)}
\label{alg:cbf}
\begin{algorithmic}[1]
\REQUIRE Robot position $\mathbf{p}$, VLA desired velocity $\vdes$,
         map obstacles $\{(\mathbf{o}_i, r_i)\}$, gain $\alpha$, safe distance $\dsafe$
\ENSURE Safe velocity command $\vcmd$
\STATE $A \leftarrow [\,]$, $b \leftarrow [\,]$
\FOR{each obstacle $i$}
  \STATE $\dsurf \leftarrow \|\mathbf{p} - \mathbf{o}_i\|_2 - r_i$
  \STATE $\hcbf \leftarrow \dsurf^2 - \dsafe^2$
  \STATE $\nabla\hcbf \leftarrow \frac{2\,\dsurf\,(\mathbf{p} - \mathbf{o}_i)}{\|\mathbf{p} - \mathbf{o}_i\|_2}$ \quad $\in \mathbb{R}^2$
  \STATE Append $\nabla\hcbf$ to $A$; append $-\alpha\,\hcbf$ to $b$
  \quad // constraint: $\nabla\hcbf \cdot \mathbf{v} \geq -\alpha\,\hcbf$
\ENDFOR
\IF{$A \cdot \vdes \geq b$ (all $i$ constraints satisfied)}
  \RETURN $\vdes$ \quad // CBF idle, no intervention
\ELSE
  \STATE $\vcmd \leftarrow \arg\min_\mathbf{v} \|\mathbf{v} - \vdes\|^2$
         subject to $A\mathbf{v} \geq b$
  \RETURN $\vcmd$ \quad // Minimum-deviation safe velocity
\ENDIF
\end{algorithmic}
\end{algorithm}

\subsection{Per-Obstacle Radii and Invisible Hazards}
Obstacle radii $\{r_i\}$ are read from the robot's on-board map.
For Isaac Sim evaluation, obstacles are rendered without visual material
(invisible to the camera), modelling map-registered hazards that no
camera-based policy can detect. The CBF continues to function because
it operates on the map, not the camera feed.

% ─────────────────────────────────────────────────────────────────────────────
\section{Experimental Setup}
\label{sec:experiments}

\subsection{Policies Evaluated}
GNM, ViNT, and NoMaD were selected because they span three major paradigms
in learned visual navigation: goal-conditioned generalist navigation, Transformer-based
foundation navigation, and diffusion-based trajectory generation.
Testing FleetSafe across these three families makes the architecture-agnostic
command-layer safety filtering claim substantially more meaningful than evaluating
minor variants of a single model class.
FleetSafe operates after the policy output at the velocity-command layer,
allowing the same CBF-QP safety filter to be applied to all three backbones
without modifying their internal design.

We evaluate all three backbones with their original weights (no fine-tuning):
\begin{itemize}
  \item \textbf{GNM} \cite{shah2023gnm}: Goal-conditioned generalist navigation
    with a conservative map-following prior. Commits to goal-directed paths.
  \item \textbf{ViNT} \cite{majumdar2023vint}: Transformer-based foundation
    navigation model with an aggressive forward waypoint prior. Highly
    goal-committed; does not back off from path obstacles.
  \item \textbf{NoMaD} \cite{sridhar2024nomad}: Diffusion-based topological
    exploration; generates trajectory distributions rather than committing
    to a single forward path. Exploration prior may naturally avoid hazards
    not on the dominant trajectory mode.
\end{itemize}

The results further reveal that invisible-hazard vulnerability is
navigation-paradigm-dependent: goal-directed policies commit to
obstacle-crossing trajectories toward the target, while diffusion-based
exploration may avoid such hazards naturally. FleetSafe intervenes selectively
only where required, preserving collision-free behaviour without degrading
naturally safe policies.

\subsection{Simulation Backends}
\paragraph{MuJoCo \cite{todorov2012mujoco}.} Procedural cylindrical obstacles in a hospital corridor.
Obstacles are visually prominent (mesh + texture). 50 seeds per
(model\,×\,scene\,×\,mode).

\paragraph{Isaac Sim \cite{nvidia2023isaac}.} Hospital corridor with 1.0\,m radius cylindrical
hazards rendered without visual material. The robot's traversal axis passes
through the obstacle, guaranteeing geometric collision in RAW mode.
50 seeds per combination. Uses RTX rendering for photoreal camera images.

\subsection{Scenes}
\begin{itemize}
  \item \textbf{Hospital Corridor}: long narrow passage with central
    obstacle(s). Primary evaluation scene.
  \item \textbf{Hospital ICU Approach}: doorframe flanking obstacles
    (off-axis). Verifies CBF does-no-harm (IR should be 0\%).
  \item \textbf{Hospital Elevator Lobby}: side-wall segment obstacles.
    Verifies CBF does-no-harm.
\end{itemize}

\subsection{PROVEN Gate}
A run is \emph{PROVEN} when all five conditions hold:
(1)~seeds\,$\geq 50$;
(2)~FleetSafe does not increase collision in any (model, scene);
(3)~FleetSafe reduces collision to $\leq 5\%$ wherever baseline exceeds $5\%$;
(4)~CBF intervention rate $>0$ in at least one scene for at least one model;
(5)~Photoreal camera confirmed (Isaac only).

% ─────────────────────────────────────────────────────────────────────────────
\section{Results}
\label{sec:results}

\subsection{Hospital Corridor: Dual-Backend Evidence}

Table~\ref{tab:main} presents the central result.
Figure~\ref{fig:fig4} shows the navigation-paradigm-dependent failure pattern
and FleetSafe's targeted intervention.

\paragraph{NoMaD invisible-hazard result.}
Unlike goal-directed GNM and ViNT (which navigate straight through the
invisible map-registered hazard toward their goal), the diffusion-based NoMaD
maintains a mean minimum distance of 1.5\,m from the hazard in RAW mode
(GNM/ViNT RAW: 0.075\,m/0.089\,m respectively --- deep into collision).
This reflects NoMaD's exploration prior: it does not commit to a fixed
forward trajectory, naturally avoiding the hazard without any safety layer.
The CBF is correspondingly idle (IR\,=\,0\%), confirming do-no-harm for NoMaD
in Isaac Sim across all 50 seeds.

The results suggest that invisible map-registered hazards are not universally
catastrophic for all visual navigation policies; rather, vulnerability depends
on the underlying navigation paradigm. Goal-directed policies commit to
obstacle-crossing trajectories toward the target, while diffusion-based
exploratory policies may avoid such hazards naturally. FleetSafe remains
effective in both regimes because the safety layer operates at the
command level rather than the representation level.

This finding identifies \emph{goal-directed navigation commitment} as the
specific structural property that makes GNM and ViNT vulnerable to invisible
map hazards. It further demonstrates that FleetSafe is \emph{selective}
rather than intrusive: safety systems are frequently criticised for
over-conservatism and unnecessary intervention that degrades nominal
performance~\cite{brunke2022safe}. Here, the same filter that prevents
100\% collision in GNM/ViNT corridor episodes adds zero intervention cost
to the already-safe NoMaD, directly addressing this concern under tested
conditions.

\paragraph{Statistical note.}
For binary collision outcomes (collision / no-collision per episode), we report
Wilson score 95\% confidence intervals. The extreme results (50/50 RAW collisions,
0/50 FS collisions for ViNT-MuJoCo and GNM/ViNT-Isaac) admit closed-form
analysis: a two-sided Fisher's exact test yields $p < 10^{-28}$ for each
comparison individually. The 95\% CI on RAW=100\% with $n=50$ is [92.9\%, 100\%]
and on FS=0\% is [0\%, 7.1\%]; the intervals are non-overlapping, confirming
significance without requiring a formal NHST at the published tier.

\begin{table}[t]
\centering
\caption{Hospital Corridor: RAW vs FleetSafe (FS). $n=50$ seeds, Wilson 95\% CI.
         Isaac NoMaD RAW=0\% reflects diffusion-based exploration (min\_dist=1.5\,m);
         goal-directed GNM/ViNT fail at 100\% (min\_dist$<$0.09\,m).}
\label{tab:main}
\begin{tabular}{llrrrr}
\toprule
Backend & Model & RAW CR (\%) & FS CR (\%) & CBF IR (\%) & Status \\
\midrule
MuJoCo & GNM   & 0.0 [0,7.1] & 0.0 [0,7.1] & 0.0         & PROVEN \\
       & ViNT  & \danger{100.0} [92.9,100] & \proven{0.0} [0,7.1] & \cbf{69.7} & PROVEN \\
       & NoMaD & 0.0 [0,7.1] & 0.0 [0,7.1] & 0.0         & PROVEN \\
\midrule
Isaac  & GNM   & \danger{100.0} [92.9,100] & \proven{0.0} [0,7.1] & \cbf{49.7} & PROVEN \\
       & ViNT  & \danger{100.0} [92.9,100] & \proven{0.0} [0,7.1] & \cbf{37.9} & PROVEN \\
       & NoMaD & 0.0 [0,7.1] & \proven{0.0} [0,7.1] & 0.0         & PROVEN \\
\bottomrule
\end{tabular}
\end{table}

\begin{figure}[htbp]
  \centering
  \includegraphics[width=\linewidth]{../figures/fig4_model_agnostic.pdf}
  \caption{Navigation-paradigm-dependent safety: FleetSafe applied to 3 VLA
           backbones across 2 physics backends. RAW (hatched) vs FleetSafe (solid).
           MuJoCo: model-specific failure (ViNT aggressive prior only).
           Isaac Sim: paradigm-specific failure --- goal-directed GNM+ViNT collide
           with invisible hazards (100\%); diffusion-based NoMaD avoids naturally (0\%).
           FleetSafe intervenes correctly in all collision-prone cases, idle otherwise.}
  \label{fig:fig4}
\end{figure}

\begin{table}[t]
\centering
\caption{Isaac Sim hospital corridor: paradigm-dependent failure and FleetSafe
         selective intervention. Min-dist is mean closest approach to hazard
         surface (r\,=\,1.0\,m cylinder). $n\,{=}\,50$ seeds, PROVEN.}
\label{tab:isaac_paradigm}
\setlength{\tabcolsep}{5pt}
\begin{tabular}{lccccl}
\toprule
Model & Paradigm & RAW CR & FS CR & CBF IR & Min-dist (RAW) \\
\midrule
GNM   & Goal-directed & \danger{100\%} & \proven{0\%} & \cbf{49.7\%} & 0.075\,m \\
ViNT  & Goal-directed & \danger{100\%} & \proven{0\%} & \cbf{37.9\%} & 0.089\,m \\
NoMaD & Diffusion/exploration & \proven{0\%} & \proven{0\%} & \textbf{0.0\%} & 1.531\,m \\
\bottomrule
\end{tabular}
\end{table}

\FloatBarrier
\subsection{CBF Intervention Rate Analysis}

Figure~\ref{fig:fig1fig2} shows collision rates and intervention rates.
In MuJoCo, the CBF intervenes 69.7\% of steps for ViNT (preventing collision)
but is idle for GNM/NoMaD (doing-no-harm). In Isaac Sim, the CBF intervenes
for goal-directed GNM (49.7\%) and ViNT (37.9\%), but remains idle for
diffusion-based NoMaD (IR\,=\,0\%), which naturally avoids the invisible hazard.
Table~\ref{tab:isaac_paradigm} makes the paradigm-dependent structure explicit.

\FloatBarrier
\subsection{Full Scene Collision Matrix}

Figure~\ref{fig:fig8} presents the complete collision rate matrix across all
model $\times$ scene $\times$ backend $\times$ mode combinations. Key observations:
(i)~Only the corridor scene with the path-intersecting obstacle produces non-zero
RAW collision; (ii)~FleetSafe reduces all corridor collisions to 0\%;
(iii)~ICU Approach and Elevator Lobby show 0\% collision in RAW mode (off-path
obstacles) and remain 0\% with FleetSafe active — confirming the filter is inert
in safe scenarios.

\begin{figure}[htbp]
  \centering
  \includegraphics[width=\linewidth]{../figures/fig8_collision_heatmap.pdf}
  \caption{Full collision rate heatmap (MuJoCo + Isaac Sim, RAW + FS modes).
           Corridor RAW cells show the failure cases; all FS cells are 0\%.
           ICU and Elevator cells confirm do-no-harm (0\% in all conditions).}
  \label{fig:fig8}
\end{figure}

\FloatBarrier
\subsection{Off-Path Scenes: Do-No-Harm Verification}

In the ICU Approach and Elevator Lobby scenes, obstacles are placed
off the robot's natural traversal axis. In all (model\,×\,backend) combinations,
RAW collision rates are 0\% and FS CBF intervention rates are also 0\%.
This confirms the CBF is quiescent in safe scenarios and does not
degrade navigation performance.

\FloatBarrier
\subsection{Safety Margin Analysis}

Figure~\ref{fig:fig5} shows the minimum obstacle distance achieved during
each episode. In the corridor scene, RAW-mode ViNT (MuJoCo) and all models
(Isaac) achieve near-zero minimum distance before collision; FleetSafe
maintains a minimum distance $\geq 0.30$\,m in all cases. Off-path scenes
show identical RAW and FS margins, confirming the filter does not perturb
safe trajectories.

\begin{figure}[htbp]
  \centering
  \includegraphics[width=\linewidth]{../figures/fig5_safety_margin.pdf}
  \caption{Minimum obstacle distance: RAW vs FleetSafe (FS). Higher is safer.
           FleetSafe maintains $\geq 0.30$\,m clearance in the corridor scene
           across all models and backends.}
  \label{fig:fig5}
\end{figure}

\FloatBarrier
\subsection{Real-Time Feasibility}

Figure~\ref{fig:fig6} shows inference latency and CBF command deviation.
GNM/ViNT achieve \SI{19}--\SI{25}{\milli\second} mean latency in Isaac Sim
(well below the \SI{100}{\milli\second} real-time threshold).
NoMaD (diffusion-based) is slower (\SI{170}--\SI{250}{\milli\second} in MuJoCo
with CPU-only inference) but remains compatible with a 5\,Hz control loop.
The \texttt{raw\_vs\_safe\_delta\_l2\_mean} metric measures the L2 norm of
the velocity correction applied by the CBF-QP. In safe scenarios this is
near zero; in the corridor scene (ViNT MuJoCo), the correction averages
\SI{0.177}{\meter\per\second}, indicating the CBF actively steers the robot
away from obstacles.

\begin{figure}[htbp]
  \centering
  \includegraphics[width=\linewidth]{../figures/fig6_latency_overhead.pdf}
  \caption{Inference latency (left two: MuJoCo; right two: Isaac Sim) and
           CBF command deviation (grey bars, right axis). GNM and ViNT operate
           well within the \SI{100}{\milli\second} real-time threshold.
           NoMaD's diffusion head is slower in CPU-only MuJoCo mode.}
  \label{fig:fig6}
\end{figure}

% ─────────────────────────────────────────────────────────────────────────────
\section{Evidence Governance}
\label{sec:governance}

We adopt a formal evidence tier ladder:
\textsc{sim-mock} $\to$ \textsc{sim-mujoco} $\to$ \textsc{sim-isaac}
$\to$ \textsc{real-proven}.
Claims may only be made at or below the tier at which evidence exists.
The PROVEN gate (§\ref{sec:experiments}) prevents premature scientific claims.
Full evidence manifests with SHA-256 hashes are available in the repository.

\paragraph{Supportable claims (v0.1).}
\begin{enumerate}
\item FleetSafe reduces collision from 100\% to 0\% for ViNT in MuJoCo
  ($n=50$, PROVEN) and for GNM/ViNT in Isaac Sim invisible hazards
  ($n=50$, PROVEN). All six cases: PROVEN.
\item Goal-directed backbones (GNM, ViNT) are vulnerable to
  camera-invisible map-registered hazards (100\% RAW collision).
  The diffusion-based NoMaD avoids naturally (0\% RAW, min\_dist=1.5\,m),
  revealing \emph{navigation-paradigm dependency} as the key risk factor.
\item CBF intervention is model-specific in MuJoCo (IR=69.7\% for ViNT,
  0\% for GNM/NoMaD) and paradigm-specific in Isaac Sim
  (IR=49.7\%/37.9\% for goal-directed GNM/ViNT; 0\% for NoMaD).
\item The filter does not degrade performance in safe scenarios:
  ICU Approach and Elevator Lobby show 0\% collision and 0\% IR in
  all 12 (model, backend) combinations tested — do-no-harm PROVEN.
\end{enumerate}

\paragraph{Not supportable from v0.1 data.}
\begin{itemize}
\item Claims about real Yahboom M3Pro hardware performance.
\item Claims about visual navigation quality (policy quality is not the
  focus; models use their original weights).
\item Claims that FleetSafe generalises to non-cylindrical obstacles,
  outdoor environments, or other robot morphologies.
\end{itemize}

% ─────────────────────────────────────────────────────────────────────────────
\section{Conclusion}
\label{sec:conclusion}

We have demonstrated that a command-layer CBF-QP filter provides
architecture-agnostic command-layer safety filtering for visual navigation
policies across multiple paradigms and simulation backends under tested
conditions. The central finding is that different embodied navigation paradigms
exhibit distinct failure structures, but the same command-layer CBF
adaptively intervenes only when required without degrading naturally safe
behaviours. This is evidenced by: (i)~collision prevention for goal-directed
GNM/ViNT in both backends; (ii)~model-specific correction for ViNT's
aggressive prior in MuJoCo; and (iii)~do-no-harm behaviour for diffusion-based
NoMaD, where IR\,=\,0\% across all 50 seeds and all three scenes.

The NoMaD result is not a failure of FleetSafe. It is evidence that the
filter is \emph{selective rather than intrusive}---a property that directly
addresses the over-conservatism critique of command-layer safety systems.

\paragraph{Future Work.}
Real-robot evidence is the next priority: a bench session on the
Yahboom M3Pro (ROS2) will validate the CBF filter under real perception
noise and actuator delay. We plan to extend the benchmark to non-cylindrical
obstacles (door frames, furniture) and dynamic agents (pedestrians).
Multi-policy ensembling and uncertainty-aware CBF gains are promising
directions for improving intervention efficiency. The paradigm-dependent
failure analysis motivates studying whether fine-tuned or hybrid
goal-directed/diffusion policies inherit the vulnerability of goal-directed
models.

\paragraph{Limitations.}
Isaac Sim results are complete for all three models (GNM, ViNT, NoMaD),
$n\,{=}\,50$ seeds per condition, PROVEN gate certified.
Real-robot evidence is pending. Claims about generalisation to outdoor
environments, non-cylindrical obstacles, or dynamic agents are not supported
by current evidence.
The surface-distance CBF assumes accurate obstacle radii are available
in the robot map; errors in radius estimation will affect the safety margin.
All results use a fixed $\alpha = 1.0$ class-$\mathcal{K}$ gain; optimal
gain tuning per environment is an open problem.
The current CBF-QP constraint uses a 1D velocity projection
($\dot{h} \approx 2d_{\text{surf}}\,(v_x \cdot \hat{d}_x)$), which is
accurate when the robot is oriented toward the obstacle but
conservative otherwise; a full 2D projection is straightforward and
planned for the next version.
All experimental evidence is from simulation; real-world sensor noise,
actuator delay, and SLAM drift are not modelled.

% ─────────────────────────────────────────────────────────────────────────────
\clearpage
\appendix
\section{Full Scene Results}
\label{app:fullscene}

Table~\ref{tab:full} reports all (model\,×\,scene\,×\,mode) results for the
MuJoCo backend (50 seeds, PROVEN). ICU Approach and Elevator Lobby scenes
place obstacles off the robot's traversal axis; all backbones achieve 0\%
collision in RAW mode, and the CBF is quiescent (IR\,=\,0\%) — confirming
the do-no-harm property.

\begin{table}[htbp]
\centering
\caption{MuJoCo full scene results — PROVEN (50 seeds). IR = CBF Intervention Rate.
         Corridor is the primary safety-evaluation scene; ICU and Elevator verify
         the do-no-harm property.}
\label{tab:full}
\setlength{\tabcolsep}{4pt}
\begin{tabular}{llcccc}
\toprule
Model & Scene & RAW CR & FS CR & CBF IR & n \\
\midrule
GNM   & Corridor       & 0.0\%  & 0.0\% & 0.0\%  & 50 \\
      & ICU Approach   & 0.0\%  & 0.0\% & 0.0\%  & 50 \\
      & Elevator Lobby & 0.0\%  & 0.0\% & 0.0\%  & 50 \\
\midrule
ViNT  & Corridor       & \danger{100.0\%}  & \proven{0.0\%} & \cbf{69.7\%} & 50 \\
      & ICU Approach   & 0.0\%  & 0.0\% & 0.0\%  & 50 \\
      & Elevator Lobby & 0.0\%  & 0.0\% & 0.0\%  & 50 \\
\midrule
NoMaD & Corridor       & 0.0\%  & 0.0\% & 0.0\%  & 50 \\
      & ICU Approach   & 0.0\%  & 0.0\% & 0.0\%  & 50 \\
      & Elevator Lobby & 0.0\%  & 0.0\% & 0.0\%  & 50 \\
\bottomrule
\end{tabular}
\end{table}

\section{Evidence Governance — PROVEN Gate Definition}
\label{app:proven_gate}

A simulation run is granted the \emph{PROVEN} status when all five conditions
hold simultaneously:

\begin{enumerate}
  \item \textbf{Seed sufficiency}: $n \geq 50$ episodes per
    (model\,×\,scene\,×\,mode) combination.
  \item \textbf{Safety monotonicity}: FleetSafe collision rate $\leq$ RAW
    collision rate for every (model, scene) pair.
  \item \textbf{Collision reduction}: In every (model, scene) where RAW
    collision rate\,$> 5\%$, FleetSafe reduces it to $\leq 5\%$.
  \item \textbf{CBF activity}: CBF intervention rate\,$> 0\%$ in at least
    one (model, scene) pair — confirming the filter is not trivially idle.
  \item \textbf{Photoreal verification} (Isaac Sim only): RTX renderer is
    enabled and confirmed active for camera image generation.
\end{enumerate}

These gates are evaluated programmatically in
\texttt{publication\_run\_scanner.py} and the result hash-committed to the
repository. No human override is permitted; a partial run (e.g., only GNM
completed) cannot be PROVEN even if individual model results pass.
The full set \{GNM, ViNT, NoMaD\} must be present.

% ─────────────────────────────────────────────────────────────────────────────
\begin{thebibliography}{99}

\bibitem{shah2023gnm}
D. Shah, B. Eysenbach, G. Kahn, N. Rhinehart, S. Levine.
\textit{GNM: A General Navigation Model to Drive Any Robot}.
ICRA 2023.

\bibitem{majumdar2023vint}
A. Majumdar, K. Yadav, S. Arnaud, Y. Ma, C. Chen, S. Silwal,
A. Jain, V. Berges, T. Sukhatme, P. Abbeel, D. Shah.
\textit{ViNT: A Foundation Model for Visual Navigation}.
CoRL 2023.

\bibitem{sridhar2024nomad}
A. Sridhar, D. Shah, C. Glossop, S. Levine.
\textit{NoMaD: Goal Masked Diffusion Policies for Navigation and Exploration}.
ICRA 2024.

\bibitem{ames2019cbf}
A. D. Ames, S. Coogan, M. Egerstedt, G. Notomista, K. Sreenath, P. Tabuada.
\textit{Control barrier functions: Theory and applications}.
ECC 2019.

\bibitem{ames2017cbf}
A. D. Ames, X. Xu, J. W. Grizzle, P. Tabuada.
\textit{Control barrier function based quadratic programs for safety critical systems}.
IEEE TAC 2017.

\bibitem{robey2021cbf}
A. Robey, H. Hu, L. Lindemann, H. Zhang, D. V. Dimarogonas, S. Tu, N. Matni.
\textit{Learning control barrier functions from expert demonstrations}.
CDC 2020.

\bibitem{emam2022cbf}
Y. Emam, P. Glotfelter, S. Wilson, A. Notomista, M. Egerstedt.
\textit{Data-driven control barrier functions for safe, multi-agent dynamics}.
IEEE LCSS 2022.

\bibitem{garcia2015survey}
J. Garcia, F. Fernandez.
\textit{A comprehensive survey on safe reinforcement learning}.
JMLR 2015.

\bibitem{alshiekh2018safe}
M. Alshiekh, R. Bloem, R. Ehlers, B. K\"onighofer, S. Niekum, U. Topcu.
\textit{Safe reinforcement learning via shielding}.
AAAI 2018.

\bibitem{achiam2017cpo}
J. Achiam, D. Held, A. Tamar, P. Abbeel.
\textit{Constrained policy optimization}.
ICML 2017.

\bibitem{falcone2011monitoring}
P. Falcone, M. Bauer, J. Svensson, J. Isgr\'o.
\textit{Predictive threat assessment via reachability analysis and set invariance theory}.
IEEE T-ITS 2011.

\bibitem{xiao2021high}
W. Xiao, C. Belta.
\textit{High relative degree control barrier functions}.
IEEE TAC 2022.

\bibitem{rauscher2016cbf}
M. Rauscher, M. Kimmel, S. Hirche.
\textit{Constrained robot control using control barrier functions}.
IROS 2016.

\bibitem{wang2017safe}
L. Wang, A. Ames, M. Egerstedt.
\textit{Safe certificate-based maneuvers for teams of quadrotors using differential flatness}.
ICRA 2017.

\bibitem{ozkil2009service}
A. G. Ozkil, Z. Fan, S. Dawids, H. Aaner, J. K. Kristensen, K. H. Christensen.
\textit{Service robots for hospitals: A case study of transportation tasks in a hospital}.
ROBIO 2009.

\bibitem{cardinaux2011video}
F. Cardinaux, D. Bhowmik, C. Abhayaratne, M. S. Hawley.
\textit{Video based technology for ambient assisted living: A review of the literature}.
J. Amb. Int. Smart Env. 2011.

\bibitem{todorov2012mujoco}
E. Todorov, T. Erez, Y. Tassa.
\textit{MuJoCo: A physics engine for model-based control}.
IROS 2012.

\bibitem{nvidia2023isaac}
NVIDIA.
\textit{Isaac Sim: A Robotics Simulation Platform}.
NVIDIA Developer Documentation, 2023.
\url{https://developer.nvidia.com/isaac-sim}

\bibitem{brockman2016openai}
G. Brockman, V. Cheung, L. Pettersson, J. Schneider, J. Schulman, J. Tang, W. Zaremba.
\textit{OpenAI Gym}.
arXiv:1606.01540, 2016.

\bibitem{macenski2022nav2}
S. Macenski, F. Mart\'in, R. White, J. G. Clavero.
\textit{The Marathon 2: A Navigation System}.
IROS 2020.

\bibitem{brunke2022safe}
L. Brunke, M. Greeff, A. W. Hall, Z. Yuan, S. Zhou, J. Panerati, A. P. Schoellig.
\textit{Safe Learning in Robotics: From Learning-Based Control to Safe Reinforcement Learning}.
Annual Review of Control, Robotics, and Autonomous Systems, 5:411--444, 2022.

\end{thebibliography}

\end{document}
